
\documentclass[10pt]{article}

% Specify the margins and text size.
\setlength{\textwidth}{6.4in}
\setlength{\textheight}{9.5in}
\setlength{\oddsidemargin}{0pt}
\setlength{\evensidemargin}{0pt}
\setlength{\topmargin}{0pt}
\setlength{\hoffset}{.05in}
\setlength{\voffset}{-.8in}

\setlength{\parskip}{5pt}
\setlength{\parindent}{0pt}

% Load some fonts and symbol packages
\usepackage{latexsym}
\usepackage{pifont}       % contains 'star' symbol for counterinsurgency handbook title
\usepackage{yfonts} 
\usepackage{amsmath}
\usepackage{amsfonts}

\usepackage{graphicx}     % actually, this is loaded by pstricks
\usepackage[T1]{fontenc}
\usepackage{ifthen}
\usepackage{pstricks,pst-grad,pst-text,pst-node,multido,pst-plot,calc,pst-3dplot}
%\usepackage[all]{xy}
%\usepackage{animate}

% The hyperref package inserts links into the pdf.
\definecolor{MyLinkColor}{rgb}{.1,.2,1}
\definecolor{MyCiteColor}{rgb}{.1,1,.2}
\definecolor{MyURLColor}{rgb}{.2,.2,.6}
\usepackage[backref=true,pagebackref=false,hyperindex,colorlinks=true,
  linkcolor=MyLinkColor,urlcolor=MyURLColor]{hyperref}


% The tweaklist package is something I found on the web.  It provides a simple interface
% for making changes to spacing used in the itemize and enumerate environments.  Comment
% this out if you don't care to use tweaklist.
\usepackage{tweaklist}
\renewcommand{\itemhook}{\setlength{\parskip}{2pt}\setlength{\parsep}%
{1pt}\setlength{\topsep}{0pt}\setlength{\itemsep}{0pt}}

\newcommand{\U}{\underline{\hspace{5pt}}}

\usepackage{listings}

\begin{document}
\pagestyle{empty}
\lstset{language=R, showspaces=false, showstringspaces=false}

\href{http://www.su.edu}{\includegraphics[height=1.75cm]{sulogo.eps}}
\vspace{-1.69cm}

{\small \hfill
\begin{tabular}{cl}
& Math 207\\& Introduction to Statistics\\
\end{tabular}
}
\setlength{\baselineskip}{1.05\baselineskip}
\medskip
\medskip

\begin{center}
\textbf{\large  First steps with the  NHANESIII Data Set}
\end{center}


\begin{itemize}
\item \textbf{Download the data set}:
\begin{itemize}
\item Create a folder called \texttt{math207} on your desktop.
\item Download the \texttt{NHANESIII.zip} file from Canvas.
\item Move the \texttt{NHANESIII.zip} file into your \texttt{math207} directory.
\item Open a terminal and type the following:\\
  \texttt{cd  Desktop}\\
  \texttt{cd  math207}\\
  \texttt{unzip  NHANESIII.zip}\\
\end{itemize}
%
\item \textbf{Open one of the files in RStudio}:
  \begin{itemize}
  \item Open RStudio.
  \item In the NNNN, type the following (where \texttt{username} is your username):\\
      \texttt{adult = read.csv("/Users/username/Desktop/math207/adult.csv", header=TRUE)}\\
      \texttt{dim(adult)}
  \item The output should be \texttt{20050  1238}  (20,050 rows and 1238 columns of data).\\
  \end{itemize}
\item \textbf{Open the code book}:
  \begin{itemize}
  \item Back in the terminal, type the following.\\
      \texttt{open  ADULT-acc.pdf}
  \end{itemize}
\end{itemize}

\vfill
\eject
\end{document}
